\subsection{Structural properties}\label{sec:struct}

The properties of Section~\ref{sec:structth} hold for the model with frozen
(explanations, certificates) or public (privacy) nuisances, and for any linear
operator; we check them for DRUP and report where they are informative.

\begin{table}[t]\centering\small
\caption{Exact explanations (Proposition~\ref{prop:affine}, Theorem~\ref{thm:cf}) on 200 random
test users: largest relative error of completeness and subset additivity
against re-propagation; rank of the top-1 item after un-logging the five
interactions with the largest attribution, five random ones, or the five most
similar ones (0 = still first); share of top-1 recommendations for which no
subset of the user's own interactions makes the second item overtake it; median
size of the minimal counterfactual explanation otherwise. Every explanation
was verified by re-scoring.}\label{tab:explain}
\begin{adjustbox}{max width=\textwidth}
\begin{tabular}{lcccc}
\toprule
Data & max.\ rel.\ error & top-1 rank after 5 removals (attr.\ / random / similarity) & no CF explanation & median CF size\\
\midrule
Coat & 6.1e-16 & 0.04 / 0.02 / 0.01 & 67\% & 0.0\\
Yahoo!\,R3 & 8.7e-16 & 1.65 / 0.27 / 1.40 & 16\% & 1.0\\
\bottomrule
\end{tabular}

\end{adjustbox}
\end{table}

\paragraph{Explanations.} Attributions recomputed against full re-propagation
are exact to floating-point precision (Table~\ref{tab:explain}), and removing the
interactions with the largest attribution moves the top item further down than
removing random or similar ones, so the attributions are faithful. The
counterfactual analysis mostly certifies that no explanation from the user's
own data exists: on Coat 94\% of the top recommendations cannot be overturned by removing any subset of the user's own interactions, on KuaiRec 52\%, and on Yahoo!\,R3 11\%; the minimal explanation otherwise has a median size of 2, 7 and 1 interactions. Such recommendations rest on population-level
evidence (through $G$), which the attribution to a user's own interactions does
not explain.

\begin{table}[t]\centering\small
\caption{Injection attack: $F$ fake users log a low-popularity target and the most
popular items. Realised hit rate of the target in real users' top-$K$, and the
fraction of real users for whom Theorem~\ref{thm:robust} certifies that the target
cannot enter the top-$K$ against \emph{any} $F$ profiles of that size (frozen
nuisances). No certificate was violated.}\label{tab:attack}
\begin{adjustbox}{max width=\textwidth}
\begin{tabular}{llcccccc}
\toprule
Data & Operator & \multicolumn{3}{c}{realised hit rate} & \multicolumn{3}{c}{certified fraction}\\
\midrule
Coat & $F=$ & 1 & 5 & 20 & 1 & 5 & 20\\
 & logged graph & 0.004 & 0.009 & 0.444 & 0.00 & 0.00 & 0.00\\
 & IPS adjacency & 0.001 & 0.269 & 0.661 & 0.15 & 0.00 & 0.00\\
 & DR adjacency & 0.000 & 0.000 & 0.029 & 1.00 & 0.94 & 0.11\\
 & DRUP & 0.000 & 0.000 & 0.028 & 1.00 & 0.93 & 0.10\\
\midrule
Yahoo!\,R3 & $F=$ & 1 & 5 & 100 & 1 & 5 & 100\\
 & logged graph & 0.065 & 0.065 & 0.199 & 0.12 & 0.00 & 0.00\\
 & IPS adjacency & 0.065 & 0.065 & 0.171 & 0.00 & 0.00 & 0.00\\
 & DR adjacency & 0.065 & 0.065 & 0.171 & 0.00 & 0.00 & 0.00\\
 & DRUP & 0.065 & 0.065 & 0.171 & 0.00 & 0.00 & 0.00\\
\midrule
KuaiRec & $F=$ & 1 & 5 & 200 & 1 & 5 & 200\\
 & logged graph & 0.002 & 0.002 & 0.288 & 0.00 & 0.00 & 0.00\\
 & IPS adjacency & 0.001 & 0.001 & 0.105 & 0.98 & 0.61 & 0.00\\
 & DR adjacency & 0.000 & 0.000 & 0.393 & 1.00 & 0.95 & 0.00\\
 & DRUP & 0.000 & 0.000 & 0.393 & 1.00 & 0.95 & 0.00\\
\bottomrule
\end{tabular}

\end{adjustbox}
\end{table}

\paragraph{Robustness.} No certificate was violated in any run
(Table~\ref{tab:attack}). With imputation-based degrees, DR adjacency and DRUP
certify every Coat user against one or two fake profiles and 93--95\% of users on
Coat and KuaiRec against five, and realised attacks with up to 100 fake users
never placed the target in a KuaiRec top-20. Certificates are informative only
for small budgets: at $F=20$ they cover 10--26\% of users, and at 200 fake users
the target reaches 39\% of KuaiRec users under DR and DRUP, more than under IPS
(10\%) and the logged graph (29\%). On Yahoo!\,R3, where DRUP selects IPS edges with
cross-fitted $W$-degrees, the fake users' degrees are attacker-controlled and no
user is certified. When the nuisances are re-fitted on the poisoned log, which
the certificates do not cover, none of the certified user--target pairs of DRUP (5{,}196 on Coat, 2{,}821 on KuaiRec)
was violated, and the realised hit rates were lower than with frozen nuisances
(Table~\ref{tab:refit} in Appendix~\ref{app:more}), because the fake users' own
imputation raises their degree and lowers their weight.

\begin{table}[t]\centering\small
\caption{nDCG of jointly private DRUP ($\delta=10^{-5}$, analytic Gaussian mechanism,
denoising rank 32 fixed in advance). Population-level nuisances come from a
public 10\% of the non-test users; ``all nuisances'' is the non-private model
with all nuisances fitted on the whole log. Other ranks in
Table~\ref{tab:privfull}.}\label{tab:privacy}
\begin{adjustbox}{max width=\textwidth}
\begin{tabular}{lccccccccc}
\toprule
Data & $\epsilon=0.5$ & 1 & 2 & 4 & 8 & 16 & $\infty$ & all nuisances & popularity\\
\midrule
Coat & 0.391 & 0.402 & 0.416 & 0.434 & 0.454 & 0.457 & 0.454 & 0.555 & 0.527\\
Yahoo!\,R3 & 0.397 & 0.437 & 0.480 & 0.518 & 0.551 & 0.574 & 0.608 & 0.652 & 0.525\\
\bottomrule
\end{tabular}

\end{adjustbox}
\end{table}

\paragraph{Privacy.} Table~\ref{tab:privacy} reports the jointly private variant of
Theorem~\ref{thm:dp}, in which only public users inform the population-level
nuisances. Useful privacy needs a large user base and a weak
guarantee. On KuaiRec, private DRUP reaches nDCG@20 of 0.596 at $\epsilon=16$ and 0.530 at
$\epsilon=8$, against 0.637 without privacy and 0.577 for popularity, and it falls below
popularity for $\epsilon\le4$. On Yahoo!\,R3 it reaches 0.574 at $\epsilon=16$ (non-private
0.652, popularity 0.525). On Coat the public-nuisance model itself is weak, because
29 public users cannot estimate the item parameters: without any noise it
reaches only 0.454, below popularity. These numbers are lower than a variant
that treats all nuisances as public would suggest, and they show the cost of an
honest privacy accounting; privatising the nuisance models with a composed
budget is left open.
